\section{ORM Performance}

\slidebackground{backgrounds/wallpaper.jpg}

\begin{frame}[fragile]{What Is ORM Performance?}
	\begin{itemize}
		\item ORM performance means getting correct results with fewer/faster database operations.
		\item Common problems:
		\begin{itemize}
			\item N+1 queries
			\item browsing huge recordsets unnecessarily
			\item searching without indexes
			\item writing in Python loops instead of batch ORM calls
		\end{itemize}
	\end{itemize}
\end{frame}

\begin{frame}[fragile]{N+1 Query Problem}
\begin{lstlisting}
# Bad: one search + one query per student
students = self.env['student.student'].search([])
for student in students:
	print(student.classroom_id.name)
\end{lstlisting}
\begin{lstlisting}
# Better: prefetch / batch read relational data
students = self.env['student.student'].search([])
students.mapped('classroom_id.name')
\end{lstlisting}
	\begin{itemize}
		\item Prefer batch operations over per-record chatter with the database.
	\end{itemize}
\end{frame}

\begin{frame}[fragile]{search vs search\_count vs search\_read}
\begin{lstlisting}
# Need records? use search
recs = self.env['student.student'].search([('active', '=', True)], limit=80)

# Need only a number? use search_count
total = self.env['student.student'].search_count([('active', '=', True)])

# Need dict rows for API/export? use search_read
rows = self.env['student.student'].search_read(
	[('active', '=', True)], ['name', 'email'], limit=80
)
\end{lstlisting}
\end{frame}

\begin{frame}[fragile]{Batch Create / Write / Unlink}
\begin{lstlisting}
# Bad
for vals in vals_list:
	self.env['student.student'].create(vals)

# Good
self.env['student.student'].create(vals_list)

# Bad
for rec in records:
	rec.write({'active': False})

# Good
records.write({'active': False})
\end{lstlisting}
\end{frame}

\begin{frame}[fragile]{Indexes and Domains}
\begin{lstlisting}
code = fields.Char(index=True)
company_id = fields.Many2one('res.company', index=True)
\end{lstlisting}
	\begin{itemize}
		\item Index fields used often in search/domain/group by.
		\item Keep domains selective (\texttt{company\_id}, \texttt{active}, dates, codes).
		\item Avoid \texttt{ilike '\%term\%'} on huge tables when exact match is possible.
	\end{itemize}
\end{frame}

\begin{frame}[fragile]{Prefetch and mapped/filtered}
\begin{lstlisting}
students = self.env['student.student'].browse(ids)
# relational prefetch helps when accessing same field on many records
classrooms = students.mapped('classroom_id')
adults = students.filtered(lambda s: s.age >= 18)
\end{lstlisting}
	\begin{itemize}
		\item Use \texttt{filtered}/\texttt{mapped} on already-loaded recordsets.
		\item Use \texttt{search} when filtering should happen in SQL.
	\end{itemize}
\end{frame}

\begin{frame}{ORM Performance Checklist}
	\begin{itemize}
		\item No N+1 loops on related fields
		\item Batch create/write/unlink
		\item Use \texttt{search\_count} for counts
		\item Limit fields with \texttt{search\_read}/\texttt{read}
		\item Index real search fields
		\item Profile before ``optimizing'' randomly
	\end{itemize}
\end{frame}
